\documentclass[12pt,a4paper]{article}

% Khai báo các gói hỗ trợ tiếng Việt, bảng biểu và toán học
\usepackage[utf8]{inputenc}
\usepackage[T5]{fontenc}
\usepackage[vietnamese]{babel}
\usepackage{geometry}
\usepackage{amsmath}
\usepackage{booktabs}
\usepackage{enumitem}
\usepackage{xcolor}
\usepackage{hyperref}

\geometry{
    a4paper,
    left=25mm,
    right=25mm,
    top=25mm,
    bottom=25mm
}

% Tùy chỉnh màu sắc đường link
\hypersetup{
    colorlinks=true,
    linkcolor=blue,
    filecolor=magenta,      
    urlcolor=cyan,
}

\title{\textbf{BIÊN BẢN NGHIỆM THU: GIAI ĐOẠN 6 -- LEAKAGE TESTS}}
\author{Dự án: How Much Intelligence Can We Get For How Little?}
\date{\today}

\begin{document}

\maketitle
\hrule
\vspace{15pt}

\section{Mục tiêu}
Giai đoạn 6 nhằm kiểm tra các invariant liên quan đến \textit{information availability}, bảo đảm rằng tại mỗi evaluation event ở thời điểm $t$, history và các thông tin đầu vào chỉ sử dụng những dữ liệu đã tồn tại trước $t$. Các kiểm tra này được xem là software correctness tests của pipeline, thay vì chỉ là các giả định được mô tả trong tài liệu.

\section{Các invariant đã kiểm tra}
Trên tập 43,027 evaluation events, các kiểm tra chính đều đạt yêu cầu.

\begin{itemize}[noitemsep]
    \item \textbf{Temporal validity:} history được xây dựng chỉ từ các interaction xảy ra trước evaluation timestamp. Kết quả được kế thừa và xác nhận từ temporal audit của Giai đoạn 4.
    \item \textbf{History nesting:} với các cửa sổ history tăng dần, history nhỏ hơn là tập con của history lớn hơn.
    \item \textbf{Full-history nesting:} mọi history hữu hạn đều là tập con của full history tương ứng.
    \item \textbf{Fixed evaluation set:} evaluation population được giữ cố định giữa các cấu hình history, không phụ thuộc vào $K$.
    \item \textbf{Fixed candidates:} candidate video được giữ nguyên giữa các cấu hình history.
    \item \textbf{Fixed labels:} nhãn \texttt{long\_view} được xác định duy nhất cho mỗi evaluation event và không phụ thuộc vào độ dài history. Trong 43,027 events, có 39,406 mẫu với \texttt{long\_view}=0 và 3,621 mẫu với \texttt{long\_view}=1.
    \item \textbf{History length consistency:} độ dài history thực tế phù hợp với các giá trị $K$ được quy định.
\end{itemize}

\section{Kiểm tra item features}
Static item metadata từ \texttt{video\_features\_basic\_1k.csv} được kiểm tra về coverage và khả năng sử dụng theo thời gian. Toàn bộ 7,388 unique evaluation videos đều có metadata tương ứng.

Các trường static metadata được xem là hợp lệ gồm: \texttt{author\_id}, \texttt{video\_type}, \texttt{upload\_dt}, \texttt{upload\_type}, \texttt{visible\_status}, \texttt{video\_duration}, \texttt{server\_width}, \texttt{server\_height}, \texttt{music\_id}, \texttt{music\_type} và \texttt{tag}.

Ngược lại, \texttt{video\_features\_statistic\_1k.csv} được loại khỏi feature candidate vì các aggregate behavioral statistics có khả năng chứa thông tin phát sinh sau thời điểm evaluation.

\section{Kết quả}
Các invariant phía data và history đã được kiểm tra đều đạt:

\begin{center}
\begin{tabular}{lc}
\toprule
\textbf{Invariant} & \textbf{Kết quả} \\
\midrule
Temporal validity & PASS \\
History nesting & PASS \\
Full-history nesting & PASS \\
Fixed evaluation set & PASS \\
Fixed candidates & PASS \\
Fixed labels & PASS \\
History length consistency & PASS \\
Static metadata coverage & PASS \\
Future aggregate statistics & EXCLUDED \\
\bottomrule
\end{tabular}
\end{center}

Một số invariant phụ thuộc vào implementation của model chưa được kiểm tra ở giai đoạn này, bao gồm \textit{constant parameter count} và \textit{frozen inference}. Hai kiểm tra này sẽ được thực hiện sau khi model và inference pipeline được triển khai.

\section{Kết luận}
Giai đoạn 6 đã hoàn tất phần kiểm định leakage ở mức data và history construction. Evaluation set, candidate, label và history được giữ độc lập với cấu hình $K$, trong khi các nguồn thông tin có nguy cơ chứa future-dependent statistics đã được loại khỏi feature candidate.

Do đó, pipeline dữ liệu hiện tại đáp ứng các invariant leakage cơ bản cần thiết để chuyển sang giai đoạn triển khai và đánh giá model.

\vspace{10pt}
\begin{center}
\large\textbf{KẾT LUẬN: GIAI ĐOẠN 6 -- HOÀN THÀNH (DATA LEVEL)}
\end{center}

\end{document}